% !TeX root = ../main.tex

\chapter{线图在非加性最短路径问题中的应用及并行化介绍}


\section{线图转换对非加性问题的适配性分析}

正如第二章所述，现实网络广泛存在如交规转弯惩罚、无人机动态油耗等路径依赖问题，导致传统的加性最短路径前提难以成立。在这些非加性（Non-additive）最短路径问题（或被称为条件权重最短路径问题）中，图网络某条边的通行代价不再是静态常量，而是受到该路径前驱历史状态约束。为了将这类具有局部马尔可夫性质的问题纳入严密的数学框架，近期研究引入了有限图的差分分析（Differential Analysis on Finite Graphs）理论，将任意路径实函数空间定义为 $H(P) = \{f : P \rightarrow \mathbb{R}\}$，并基于特定路径的差分算子 $\partial f(\alpha | p)$ 实现了对各类复杂代价函数的严格分类\cite{durand2024generalized}。

依据该微分形式理论，任何复杂路径的代价都可以被精确拆解为其局部差分的叠加：
\begin{equation}
    f(p) = \sum_{\alpha \in p} \partial f(\alpha | p)
\end{equation}
其中，$P$ 表示图 $G$ 中所有合法路径的集合，$p \in P$ 为某一条具体路径；$\alpha \in E_G$ 表示构成该路径的某一条有向边；算子 $\partial f(\alpha | p)$ 定义为在给定前置历史路径段的情况下，拓展边 $\alpha$ 接入该路径所带来的局部代价增量（即，$\partial f(\alpha|p) = f(p \cup \{\alpha\}) - f(p)$）。

在此基础上，若到达某条目标边 $\alpha$ 的差分代价变动 $\partial f$ 仅受其紧邻的前面 $k$ 条边（即有限深度的历史步骤）的条件影响，则该代价函数在拓扑学上被称为 \textbf{$k$-加性（$k$-additive）函数}，记作 $f \in F^{(k)}_G$。在这一标准下，传统可直接求解的纯加性代价函数被统一归类为 0-加性函数（$F^{(0)}_G$），而大多数涉及转弯角度或前驱关联约束的实际物理模型，均可被抽象降级为条件受限的 1-加性或 2-加性网络路线问题。

针对这种突破了零阶马尔可夫性的 $k$-加性前驱依赖问题，\textbf{线图（Line Graph）转换理论}展现出了良好的结构适配性与算法应用价值：
\begin{itemize}
    \item \textbf{网络状态的降维解耦：} 对于具有 1-加性复杂代价的原图 $G(V_G, E_G)$，其前驱约束使得节点间的状态过渡相互关联。通过将其同态映射为线图 $L(G)$，原图中的有向边映射为线图中的实体节点（即 $V_{L(G)} = E_G$）。这种拓扑转换将原图中附带领航方向的动态历史轨迹转化为线图中静态的独立节点。
    \item \textbf{非加性约束的静态显式加权：} 在原图 $G$ 中，受历史条件所限的转移代价 $\partial f(\beta | \alpha)$（从前驱边 $\alpha$ 跨越节点转移到边 $\beta$ 所产生的额外代价），被等效地合并到了线图 $L(G)$ 中新生成的相邻节点对之间的初始静态边权重中。
\end{itemize}

综上所述，通过线图转换这一空间拓扑工具，我们能够在保留完整物理约束信息的条件下，将源网络上受历史制约的 $k$-加性寻路问题，在数学映射层面等效地重构为替代图 $L(G)$ 上的 0-加性基准问题。这种拓扑变换为全局寻优算法的设计提供了坚实的理论基础。基于此，Durand 等人\cite{durand2024generalized}提出了条件权重最短路径（Conditional Weighting Shortest Path, CWSP）算法框架。下文将首先对该经典算法的具体设计与等价性证明进行详细论述，并为后续提出该算法的并行化改进方案作铺垫。

\section{基于线图的算法设计}

基于前节理论中代价函数的微分分解与拓展思维，现有研究已经建立了一套完整的条件权重最短路径（CWSP）算法框架\cite{durand2024generalized}，用于求解包含复杂前驱约束的最优路径问题。本节将对该框架进行深入剖析，其核心思路分为拓扑转换、条件加权映射与替代图寻路三个层序步骤。

\begin{figure}[H]
    \centering
    \begin{tikzpicture}[
        node distance=1.5cm and 2cm,
        box/.style={rectangle, draw, rounded corners, minimum width=3.5cm, minimum height=1.5cm, align=center, thick, fill=blue!10},
        arrow/.style={-stealth, thick},
        label/.style={font=\small, text width=3cm, align=center}
    ]
        
        % Nodes
        \node[box, fill=gray!10] (input) {原图网络 \\ $G=(V_G,E_G)$ \\ (带前驱物理约束)};
        \node[box, below=of input] (phase1) {第一阶段: 拓扑转换 \\ 节点化连边 \\ $V_H \leftarrow E_G, \, E_H \leftarrow (\alpha,\beta)$};
        \node[box, below=of phase1] (phase2) {第二阶段: 静态加权 \\ 差分增量映射 \\ $W_H \leftarrow \partial f(\beta|\alpha)$};
        \node[box, below=of phase2] (phase3) {第三阶段: 路径寻优 \\ 标准Dijkstra求解 \\ $p^* \leftarrow \min W_H(q)$};
        \node[box, fill=green!10, below=of phase3] (output) {全局最优解 \\ $p^* \in P_G(s,t)$};
        
        % Arrows
        \draw[arrow] (input) -- (phase1);
        \draw[arrow] (phase1) -- (phase2) node[midway, right, label] {无权高阶线图 \\ $L^{[k]}(G)$};
        \draw[arrow] (phase2) -- (phase3) node[midway, right, label] {带权0-加性替代图 \\ $H(V_H,E_H,W_H)$};
        \draw[arrow] (phase3) -- (output);
        
        % Decorative box framing the algorithm phases
        \draw[dashed, thick, gray] 
            ([xshift=-3cm, yshift=0.5cm]phase1.north west) rectangle 
            ([xshift=3.5cm, yshift=-0.5cm]phase3.south east);
        \node[anchor=south west, text=gray, font=\bfseries] at ([xshift=-3cm, yshift=0.5cm]phase1.north west) {CWSP算法核心步骤};
            
    \end{tikzpicture}
    \caption{非加性约束下的线图转换与最短路径求解（CWSP框架）流程}
    \label{fig:cwsp_pipeline}
\end{figure}

\subsection{第一阶段：线图的拓扑替代构造}
算法的第一步旨在将原图 $G=(V_G, E_G)$ 在连边上的隐式约束，物理显式化为图节点的邻接问题。该过程定义了原图到其替代图（Alternative Graph，本质即为线图 $H = L(G)$）的映射逻辑：
\begin{enumerate}
    \item \textbf{节点生成：} 为原图 $G$ 中的每一条有向边 $\alpha \in E_G$ 在新图 $H$ 中创建一个对应的单一实体节点。即新图节点集由原图边集构成：$V_H = \{ u \mid u = \alpha \in E_G \}$。
    \item \textbf{连边构建：} 若在原图 $G$ 中边缘 $\alpha$ 和有向边 $\beta$ 能够顺畅过渡（即边 $\alpha$ 的终点与边 $\beta$ 的起点为同一节点），则在线图 $H$ 中创建一条连接前驱节点与后续节点的新边，即 $E_H = \{ (\alpha, \beta) \mid \alpha, \beta \in V_H, \text{且~} \alpha_2 = \beta_1 \}$。
\end{enumerate}

综上拓扑映射规则，该阶段具体的构造过程如算法 \ref{alg:line_graph} 所示。

\begin{algorithm}[H]
\caption{线图的拓扑替代构造}\label{alg:line_graph}
\KwIn{原图 $G = (V_G, E_G)$}
\KwOut{线图 $L(G)$}
$V_{L(G)} \leftarrow \{(\alpha) \mid \alpha \in E_G\}$\;
$E_{L(G)} \leftarrow \emptyset$\;
\ForEach{$\alpha \in E_G$}{
    \ForEach{$\beta \in E_G$}{
        \If{$\alpha_2 = \beta_1$}{
            $E_{L(G)} \leftarrow E_{L(G)} \cup \{(\alpha, \beta)\}$\;
        }
    }
}
$L(G) \leftarrow (V_{L(G)}, E_{L(G)})$\;
\Return $L(G)$\;
\end{algorithm}

如算法 \ref{alg:line_graph} 所示，该生成过程的核心依赖于遍历原图的边集合 $E_G$。第 4-5 行完成了替代图 $L(G)$ 的节点集和边集的初始化。第 6-12 行通过双层循环枚举所有可能的原图有向边对 $(\alpha, \beta)$，并通过判断前驱边的物理终点 $\alpha_2$ 是否等同于后继边的物理起点 $\beta_1$ 来决定是否在替代图中建立新的有向连接。这种基于有向边前后继关系的重组，构成了替代图 $H$ 的整体拓扑骨架。

\subsection{第二阶段：非加性约束的静态加权}
替代图 $H$ 生成后，需要将蕴含历史约束的非加性网络代价内化到替代图的全新边集合 $E_H$ 中。假定特定问题的非加性网络最多取决于 $k$ 级历史状态，对于生成的替代图上的每条转移连边 $(\alpha, \beta)$，运用公式（3.1）中引入的偏导数评价算子即可算出其固化的静态网络权重：
\begin{equation}
    W_H(\alpha \to \beta) = \partial f(\beta \mid \alpha)
\end{equation}
此时，无论原图问题具有何种形式的约束，都已经通过局部差分值转化为替代图 $H$ 中恒定且静态的网络权重。至此，含有条件约束的复杂网络，已被等效转化为标准权重的纯加性有向图。这一阶段权重映射的整体逻辑如算法 \ref{alg:conditional_weight} 所示。

\begin{algorithm}[H]
\caption{替代图（线图）的条件加权}\label{alg:conditional_weight}
\KwIn{原图 $G = (V_G, E_G)$，复杂代价函数 $f$，代价函数的加性阶数 $k$}
\KwOut{带权替代图 $H = (V_H, E_H, W_H)$，其中 $H$ 本质为高阶线图 $L^{[k]}(G)$}
$H \leftarrow L^{[k]}(G) $ \tcc*{利用算法 \ref{alg:line_graph} 构造高阶线图}
\ForEach{$(u_H, v_H) \in E_H$}{
    $(u_1, \dots, u_{k+1}) \leftarrow u_H$\tcc*{其中 $l(u_H)=l(v_H)=k+1$}
    $(v_1, \dots, v_{k+1}) \leftarrow v_H$\;
    $\beta \leftarrow (u_{k+1}, v_{k+1})$ \tcc*{提取最新扩展边 $\beta \in E_G$}
    $w_{(u_H, v_H)} \leftarrow \partial f(\beta \mid u_H)$ \tcc*{计算差分并赋静态权重}
}
$V_H \leftarrow V_H \cup V_G$\;
\ForEach{$u_H \in V_H$}{
    $u \leftarrow (u_H)_1$\;
    $\gamma \leftarrow (u, u_H)$ \tcc*{从原图节点进入替代图 $H$ 的连边}
    $E_H \leftarrow E_H \cup \{\gamma\}$\;
    $w_\gamma \leftarrow f((u_H)_1, \dots, (u_H)_k)$\;
}
$W_H \leftarrow \{w_e \mid e \in E_H\}$\;
\Return $H = (V_H, E_H, W_H)$\;
\end{algorithm}

算法 \ref{alg:conditional_weight} 详细展示了将动态非加性代价固化为静态权重的计算过程。算法第 1 行首先调用前述逻辑构造高阶替代线图结构。第 2-6 行依次对替代图中的每一条新生成连边 $(p, q)$ 进行遍历，捕捉发生状态过渡的最新连边 $\beta$，随后应用偏导数算子 $\partial f$ 求取对应的局部差分代价，以此作为该条连边的固有静态权重。此外，算法在第 8-13 行针对由原网络起起始点寻路并过渡到替代图的“入界连接边”，单独计算了其初始状态下的边界代价，以此保证算法寻优在总体尺度上的严格数学对齐。

\subsection{第三阶段：基于等效替代图的路径寻优}
完成图结构的拓扑替代及权重重分配后，替代图 $H$ 已转化为标准的 0-加性有向图，此时可直接接入经典的 Dijkstra 算法在此新网络上进行全局寻优。由于模型物理意义的改变，应用 Dijkstra 算法时只需对其停止条件进行相应的逻辑更新。

在原始最短路径问题中，寻路过程当目标节点 $t \in V_G$ 被从优先队列弹出时即告结束。但在替代图 $H$ 结构下，节点实质为原图空间中的一条“转移线段”。因此，Dijkstra 算法的停止条件转变为：当优先队列弹出的当前节点（即某条有向边）的终点，恰好指向原问题的全局目标节点 $t$ 时，寻路探索即告完成。随后，运用回溯法依据前驱指针（predecessor）即可输出映射回原图的完整最小代价通路 $p^*$。针对转换后图结构的单源最短路寻优完整步骤如算法 \ref{alg:dijkstra_sub} 所示。

\begin{algorithm}[H]
\caption{基于替代图（线图）的 Dijkstra 寻路算法}\label{alg:dijkstra_sub}
\KwIn{带权替代图 $H = (V_H, E_H, W_H)$，起点 $s$，终点 $t$}
\KwOut{最优路径 $p^*$}
\ForEach{$v_H \in V_H$}{
    $\text{cost}[v_H] \leftarrow \infty$\;
    $\text{predecessor}[v_H] \leftarrow \emptyset$\;
}
\tcc{此处省略了与起点相关的初始化细节，默认包含从起点关联起步的过程}
$\text{cost}[s'] \leftarrow 0$ \tcc*{$s'$表征线图层级下对应源节点的特定起始状态}
$Q \leftarrow V_H$\;
\While{$Q \neq \emptyset$}{
    $u_H \leftarrow \text{extract-min}(Q)$\;
    \ForEach{连边 $\alpha_H = (u_H, v_H) \in E_H$}{
        \If{$\text{cost}[v_H] > \text{cost}[u_H] + w_{\alpha_H}$}{
            $\text{cost}[v_H] \leftarrow \text{cost}[u_H] + w_{\alpha_H}$\;
            $\text{predecessor}[v_H] \leftarrow u_H$\;
        }
        \If{$(v_H)_{k+1} = t$}{
            $u_{\text{last}} \leftarrow v_H$\;
            \textbf{break}\;
        }
    }
}
$p^* \leftarrow (t)$\;
$u_H \leftarrow u_{\text{last}}$\;
\While{$l(u_H) = k+1$}{
    $v_H \leftarrow u_H$\;
    $u_H \leftarrow \text{predecessor}[v_H]$\;
    $p^* \leftarrow ((u_H)_{k+1}, p^*)$\;
}
$p^* \leftarrow ((v_H)_1, \dots, (v_H)_k, p^*)$\;
\Return $p^*$\;
\end{algorithm}

结合算法 \ref{alg:dijkstra_sub} 可知，在寻路探索阶段的前 14 行，操作逻辑与标准 Dijkstra 算法保持高度一致。主要变化出现在第 15 行为目标节点的核验机制：由于替代图 $H$ 中遍历的节点本质已升级为原图的行进边集，故寻路停止标志更新为针对当前搜寻前沿的终点状态 $v_{k+1}$ 执行判定。一旦该终点触及全局目标节点 $t$，循环即可中止。并在算法第 20-25 行，通过向后递推追踪前驱指针，将替代图域中的节点迹线反向解耦拼装，最终呈现并返回实际地理平面上的单源最优路径。

\subsection{算法理论最优性的论证}
为进一步巩固上述算法设计的理论正确性，现对该 CWSP 算法所求得的最短路径等价于原图 $G$ 非加性问题的最优解，作如下严密证明：

\textbf{定理 3.1（最优性等价定理）} \textit{设给定有向图 $G(V_G,E_G)$ 及由源点 $s$ 至终点 $t$ 的 $1$-加性路径寻优问题，目标函数为 $\min_{p \in P_G(s,t)} f(p)$。设 $H(V_H,E_H,W_H)$ 为依据上述 CWSP 算法第一、二阶段条件加权构造出的线图，则在该图中依据第三阶段求得的纯加性最短路径解集 $\arg\min_{q \in P_H} \sum_{e \in q} W_H(e)$，其必然映射回 $P_G(s,t)$ 并对应 $f(p)$ 的全局最优解。}

\textbf{证明:}

\textbf{步骤一：解空间的双射映射性质（Bijective Mapping of Solution Space）}  
设 $p^* \in P_G(s,t)$ 为原图中的一条有效路径边序列，记为 $p^* = \langle \alpha_1, \alpha_2, \dots, \alpha_n \rangle$，起点判定条件为 $(\alpha_1)_1 = s$，终点为 $(\alpha_n)_2 = t$。依据 CWSP 第一阶段生成规则 $V_H = E_G$ 且 $E_H = \{ (\alpha_i, \alpha_{i+1}) \}$，该序列 $p^*$ 的邻接特征在 $H$ 中唯一显化为一条有效节点路径 $q = \langle v_1, v_2, \dots, v_n \rangle \in P_H$，满足 $v_i = \alpha_i \in V_H$。这也论证了 $P_G$ 与 $P_H$ 解空间之间存在严格的一一双向映射算子 $\Phi : p^* \mapsto q$。

\textbf{步骤二：非加性向加性的代价等价性（Cost Equivalence）}  
基于代价函数的差分分解定义与 $1$-加性截断边界属性，原图中路径 $p^*$ 的组合代价可递推展开：
\begin{align}
    f(p^*) &= f(\langle \alpha_1, \alpha_2, \dots, \alpha_n \rangle) \notag \\
           &= f(\alpha_1) + \sum_{i=1}^{n-1} \partial f(\alpha_{i+1} \mid \langle \alpha_1, \dots, \alpha_i \rangle) \notag \\
           &= f(\alpha_1) + \sum_{i=1}^{n-1} \partial f(\alpha_{i+1} \mid \alpha_i)
\end{align}
而对应由 CWSP 第二阶段权值定义的映射线图路径 $q$ ，其 $0$-加性累积模型为：
\begin{equation}
    W_H(q) = w_{enter}(\alpha_1) + \sum_{i=1}^{n-1} W_H(v_i, v_{i+1})
\end{equation}
定义起始惩罚 $w_{enter}(\alpha_1) = f(\alpha_1)$ 及连边静态权值项 $W_H(v_i, v_{i+1}) = \partial f(\alpha_{i+1} \mid \alpha_i)$ 后，上述两式可建立如下等价关系：$\forall p^* \in P_G(s,t), \, f(p^*) \equiv W_H(\Phi(p^*))$。

\textbf{步骤三：最优解一致性（Optimality Convergence）}  
设经由 CWSP 第三阶段寻路，在 $H$ 平面上求得 $\min_{q \in P_H} W_H(q)$ 对应的最短线路 $q^*$。假定在原图 $G$ 中存在一条总代价更小的有效路径 $p'$ 满足 $f(p') < f(\Phi^{-1}(q^*))$；由于定理证明的映射等价性 $f(p^*) \equiv W_H(q^*)$ 约束，必然在线图 $H$ 中也存在一条对应的有效同构轨迹 $q' = \Phi(p')$，满足 $W_H(q') < W_H(q^*)$。这与“$q^*$ 是 $H$ 中的最小权重路径”这一前提假设相矛盾。
因此，假设不成立，$f(\Phi^{-1}(q^*))$ 必然是全局绝对下确界。这表明了 CWSP 算法解得的 $q^*$ 完全等价于原问题网络的最优解。定理论证闭环，证毕。

需要指出的是，这种在空间维度上扩展网络的 CWSP 算法机制，在线图模型拓展时将面临边数显著增加的情况：即 $|E_{L^m(G)}| \leq n^{m+2}$（$n$为节点数）。这种节点规模的几何级数增长虽然降低了逻辑判断的复杂度，但对传统单线程 CPU 的计算能力带来了特定的挑战。针对这类大规图层遍历与计算需求，引入基于 GPU 底层的 CUDA 高并发阵列进行并行运算，成为了一种有效的解决途径。


\section{算法的并行优化}

针对 CWSP 算法在线图转换及寻路过程中面临的空间规模扩张与计算瓶颈，本节基于 GPU 的 CUDA 并行架构，对算法的三个核心阶段提出针对性的并行化改进方案。主要改进策略包括线图拓扑结构的并行构建、条件权重的并行计算，以及采用适配 GPU 架构的并行最短路径算法替代传统的 Dijkstra 算法。

\begin{figure}[H]
    \centering
    \begin{tikzpicture}[
        node distance=0.8cm and 1cm,
        stepbox/.style={rectangle, draw, rounded corners, minimum width=6cm, minimum height=1.2cm, align=center, thick, fill=orange!10},
        arrow/.style={-stealth, thick},
        font=\small
    ]

    % Draw Grid for GPU
    \draw[thick, draw=green!60!black, rounded corners, loosely dashed] (-4.5, 1.5) rectangle (4.5, -6.5);
    \node[anchor=north west, font=\bfseries\large, text=green!50!black] at (-4.2, 1.2) {GPU (CUDA) 并行架构化};

    % Phase 1: Line Graph Construction
    \node[stepbox] at (0, -0.5) (csr) {第一阶段: 拓扑高并发构建 \\ (CSR度数扫描与前缀和无锁映射)};
    
    % Phase 2: Conditional Weight SIMT
    \node[stepbox] at (0, -3.0) (simt) {第二阶段: 条件权重提取 \\ (SIMT多分支协同与并发访存优化)};
    
    % Phase 3: $\Delta$-Stepping
    \node[stepbox] at (0, .5.5) (delta) {第三阶段: $\Delta$-Stepping寻优 \\ (无序波前松弛与共享内存并发探底)};

    \draw[arrow] (csr) -- (simt);
    \draw[arrow] (simt) -- (delta);

    \end{tikzpicture}
    \caption{基于CUDA架构的线图寻路并行优化策略实施框架}
    \label{fig:parallel_optimization}
\end{figure}

\subsection{线图拓扑结构的并行构建}

在线图生成阶段，原图 $G$ 中每条有向边的拓展及其相邻关系的判定过程相互独立，不存在数据依赖性，这为细粒度的数据并行提供了可能。

在具体并行实现中，必须摒弃传统串行算法中基于动态内存分配的链表或双重循环遍历。为适配 GPU 的内存分层架构，本文采用高聚合的\textbf{压缩稀疏行（Compressed Sparse Row, CSR）}数据格式将原图结构存储于设备全局内存（Global Memory）中，并采用以“边”为粒度的两阶段并行映射策略。

第一阶段为\textbf{拓扑度数扫描与空间预分配}。由于原图中各条边引出的后继边数量（即节点的出度）参差不齐，为了防止各线程在写入连续显存时发生写冲突（Race Condition），首先分配一维的 CUDA 线程网格与线程块。每个线程（Thread）独立读取原图的一条连边 $\alpha$，统计其尾节点对应的合法相连后继边 $\beta_i$ 的数量（即新生成的替代图出度）。随后，调用高效的并行前缀和（Exclusive Prefix Sum / Scan）算法，对所有线程生成的度数数组进行扫描，由此算出替代图总边数并为其预分配连续的稀疏存储空间。同时，该前缀和的结果等同于为每条原图边分配了在全局写入数组中的绝对偏移量（Offset）。

第二阶段为\textbf{无锁化拓扑实例化}。重新启动并行的映射核函数（Kernel），各计算线程根据第一阶段获取的独占偏移量索引，将合法的 $(\alpha, \beta_i)$ 状态转移关系直接、独立地写入 CSR 结构的目标地址中。这种基于前缀和配合 CSR 格式的并行生成模式，实现了完全无锁（Lock-free）的高并发内存加载，极大缩减了大量动态构建替代图拓扑由于互斥锁导致的等待过载。

\subsection{条件权重的并行计算}

线图结构的重构完成后，所有新生成的转移关联边及其对应的原图映射信息已全部存储于连续数组中。此时，为线图赋权的偏导数算子 $\partial f(\beta \mid \alpha)$ 的计算是一个典型的单指令多数据（SIMT）操作。

在这一优化步骤中，直接针对海量新生成的替代图边集合 $E_H$ 启动一维计算核函数。每个执行线程依据其全域线程索引（Global Thread ID），精确映射并提取出一条拓扑转移变 $e_H = (\alpha, \beta)$。随后，线程依据内存中的几何与前趋历史状态属性，并行地对每一过渡关系（如计算转弯角度惩罚、油耗受阻损耗等）进行评价，继而产生代价差分增量并将其固化为静态图权值。

为进一步榨取计算管线性能，此阶段必须引入以下两项底层显存与并发约束策略：首先是\textbf{合并内存访问（Memory Coalescing）}。由于前阶段所建立的替代图连边已由连续的 CSR 格式紧凑存储于主显存，相邻计算线程在提取 $(\alpha, \beta)$ 指针特征时，其内存寻址呈现出高度的空间局部性，能够以宽总线事务的数据块形式并发出队，从硬件层面填满极速的显存带宽吞吐；其次是\textbf{线程束控制冲突抑制（Warp Divergence Control）}。复杂的物理重定权往往依赖“是否合规转弯”、“死路惩罚阈值”等多分支的 \texttt{if-else} 逻辑评判。此并行设计借由条件算子与内部谓词判定（Predication）等形式有效替换显式的控制流拆分，借以保证指派往同一线程束（Warp，即在 NVIDIA 构架中 32 个协同处理的子线程阵列）内的基础指令步骤绝对对齐。通过这些建立在微架构机制之上的严格规约，非加性权重提取分析的耗时将被极限压缩在局部内存访问常数级别。

\subsection{GPU友好的并行寻优算法替代}

在完成线图重构与赋权后，原有的 CWSP 框架依赖标准 Dijkstra 算法进行 0-加性寻路。然而，传统 Dijkstra 算法的核心在于维护一个严格的全局优先队列（Priority Queue），在每一步都必须严格取出当前距离初始节点最近的那个未访问顶点来执行松弛更新。这种强顺序依赖性（Sequential Dependency）导致它极难在单指令多线程的 GPU 环境上获得良好的并行扩展性。

为了突破这一极端的顺序依赖算力瓶颈，本文采用适配粗粒度数据分布的 $\Delta$-Stepping 算法来替代原始的 Dijkstra 同步模型框架。这是一种高效糅合了 Dijkstra 严密正确性以及 Bellman-Ford 算法无界并行松弛优势的动态折中方案。应用针对于 GPU 加速改造的核心改良逻辑如下：
\begin{enumerate}
    \item \textbf{并行松弛控制与共享内存结构映射：} $\Delta$-Stepping 引入“桶（Bucket）”数组用于彻底取代树型优先队列。该机制设定步长距离参量 $\Delta$，将图中尚未求得最终最优值的距离态划分至相互邻接的操作分区。其中任一桶 $B_i$ 内收纳所有起始路径探求边距正落于区段 $[i\Delta, (i+1)\Delta)$ 之间的替代节点的集合。在并发调度上，这部分局部活跃数据会被优先转储至靠近处理中心、访问延迟极低的每个线程块 \textbf{共享内存（Shared Memory）}内。
    \item \textbf{无序并发探测限制与原子性竞夺处理：} 针对于汇聚于同个操作桶区内的波前（Frontier）节点序列，因其具有相逼的短径延展力，GPU 内核准许大规模开启并行处理而不必维持队列内外的精确次序。然而在此并行松弛阶段中，数以百计的工作线程极易出现向同一个未知目的地同时递发多条不同探测值（不同“轻边”的寻路尝试）的问题。为应对这种极端的数据争用，核函数内部强制植入了\textbf{原子性探底更新方法（Atomic Minimum）}，如调用 \texttt{atomicMin} 指令保障并行松弛探寻仅保留真正的成本全局最优解落库。
    \item \textbf{全局屏障同步（Barrier Synchronization）与循环协同：} 桶内被激活探索的前沿往往会导致同个区间端内产生高频的复发唤醒。由于这些邻接跃迁可能交替产生依赖，在实现内外层循环转换的调度中需按批次布置明确的内核阻断设备指令（如利用新一轮内核启动来触发 \texttt{cudaDeviceSynchronize()} 同步点），以此杜绝在迭代清理完桶内轻度权重依赖或者处理大步跃迁边界重边（延展代偿 $>\Delta$）的过程中导致的全局逻辑重越穿透乱序故障。最终构建了从桶内“并行猛攻”到下域“轮转迁移”的高能效波纹扩散寻优面。
\end{enumerate}

综上改进，通过在线图中引入基于 GPU 的 $\Delta$-Stepping 短路径求解方案，本文不仅彻底规避了优先队列并发处理中的锁争用开销，还在保证全局最优解精度的前提下，充分利用了线图扩展后的巨量数据连通性特征，取得了更为显著的算法加速比。

\section{本章小结}

本章聚焦于非加性最短路径问题的求解机制与底层加速瓶颈，详细论述了图拓扑转换与现代 GPU 并行计算架构相结合的重构求解框架。

首先，针对交通网络等场景中广泛存在的历史前驱路径依赖痛点，本章引入了有限图差分分析与 $k$-加性函数理论，深层次剖析了线图结构处理非加性约束的天然适配性。通过拓扑空间转换，原图中受历史状态制约的动态过渡代价被精确解耦，并降维映射为替代图中的常数级静态边权重，从而将繁杂的非加性寻路问题无损转化为标准的 0-加性基准问题。在此数学框架下，本章详细推演了具有极高代表性的条件权重最短路径（CWSP）算法的执行体系，并在此基础上提供了严密的定理证明，保证了线图映射解与原寻优空间全局最优解之间的绝对双射等价性。

其次，针对引入线图模型后图网络规模必定发生几何爆炸而引发的传统串行算力灾难，本章突破性地提出了一整套贴合 GPU 硬件微架构的图算法深度并行化优化策略。在拓扑构建层，利用无锁化的并行前缀和协同压缩稀疏行（CSR）数据格式，彻底规避了动态显存分配带来的竞态访问锁死；在条件赋权层，依托内存合并访问与强制的线程束（Warp）分化抑制规则，实现了高吞吐的单指令多数据（SIMT）并发评价；而在最后也是致受限最深的全局寻路松弛层，本章成功引入并改良了松散依赖的 $\Delta$-Stepping 算法，借由共享内存转储、桶（Bucket）级迭代与原子操作（AtomicMin）的深度调度，彻底消融了传统 Dijkstra 中因优先队列高度串行而引起的并发死锁。这一系列立足于底层体系结构的核心再造工程，为下一章展开真实世界的大规模路网实验评估及框架部署奠定了坚实的工程基石。
